\lecture{16}{\texorpdfstring{Экспоненциальная\\ устойчивость}{Экспоненциальная устойчивость}}

\sourcevideo
    {https://www.youtube.com/playlist?list=PLOzRYVm0a65fhKDS8cYIC7YCuVGJ4FOJx}
    {Week 5: Lecture 16: Exponential Stability}

Неравенство
\[
    \dot V\le-\gamma V
\]
гарантирует экспоненциальное убывание функции Ляпунова, но для
экспоненциальной устойчивости состояния нужно дополнительно сравнить
\(V(x)\) с расстоянием до равновесия. В этой лекции такой переход
выполняется для градиентного потока сильно выпуклой функции.

\section{Квадратичный критерий}

Рассмотрим автономную систему
\[
    \dot x=g(x)
\]
с равновесием \(x^\star\), то есть \(g(x^\star)=0\), и обозначим
\[
    e=x-x^\star.
\]

\begin{theorem}[Критерий экспоненциальной устойчивости]
Пусть существует непрерывно дифференцируемая функция
\(V\colon\RR^n\to\RR\) и положительные константы
\(c_1,c_2,c_3\), для которых
\[
    c_1\lVert e\rVert^2
    \le
    V(x)
    \le
    c_2\lVert e\rVert^2
\]
и
\[
    \dot V(x)
    \le
    -c_3\lVert e\rVert^2.
\]
Тогда
\[
    \lVert x(t)-x^\star\rVert
    \le
    \sqrt{\frac{c_2}{c_1}}
    \exp\left(
        -\frac{c_3}{2c_2}(t-t_0)
    \right)
    \lVert x(t_0)-x^\star\rVert.
\]
\end{theorem}

\begin{proof}
Из верхней оценки на \(V\) следует
\[
    \lVert e\rVert^2\ge\frac1{c_2}V(x),
\]
поэтому
\[
    \dot V\le-\frac{c_3}{c_2}V.
\]
Интегрирование даёт
\[
    V(x(t))
    \le
    \exp\left(
        -\frac{c_3}{c_2}(t-t_0)
    \right)
    V(x(t_0)).
\]
Используя обе квадратичные оценки, получаем
\[
\begin{aligned}
    c_1\lVert e(t)\rVert^2
    &\le
    V(x(t))
    \\
    &\le
    c_2
    \exp\left(
        -\frac{c_3}{c_2}(t-t_0)
    \right)
    \lVert e(t_0)\rVert^2.
\end{aligned}
\]
После извлечения квадратного корня получается требуемая оценка.
\end{proof}

Таким образом, в стандартной записи
\[
    \lVert e(t)\rVert
    \le
    M e^{-\alpha(t-t_0)}\lVert e(t_0)\rVert
\]
можно взять
\[
    M=\sqrt{\frac{c_2}{c_1}},
    \qquad
    \alpha=\frac{c_3}{2c_2}.
\]
Знак минус в условии \(\dot V\le-\gamma V\) принципиален: без него
функция Ляпунова может возрастать.

\section{Сильно выпуклый градиентный поток}

Рассмотрим задачу
\[
    \min_{x\in\RR^n}f(x)
\]
и градиентный поток
\[
    \dot x=-\nabla f(x).
\]
Пусть функция \(f\) является \(\mu\)-сильно выпуклой и
\(L\)-гладкой, где
\[
    0<\mu\le L.
\]
Тогда минимум \(x^\star\) единственен и
\[
    \nabla f(x^\star)=0.
\]

Из сильной выпуклости при выборе базовой точки \(x^\star\) следует
\[
    f(x)-f^\star
    \ge
    \frac\mu2\lVert x-x^\star\rVert^2.
\]
Из леммы об убывании и \(L\)-гладкости получаем
\[
    f(x)-f^\star
    \le
    \frac L2\lVert x-x^\star\rVert^2.
\]
Следовательно,
\[
    \frac\mu2\lVert x-x^\star\rVert^2
    \le
    f(x)-f^\star
    \le
    \frac L2\lVert x-x^\star\rVert^2.
\]
Именно эта двусторонняя оценка позволяет использовать ошибку по
значению функции как квадратично сравнимую функцию Ляпунова.

Сильная выпуклость также влечёт PL-неравенство
\[
    \frac12\lVert\nabla f(x)\rVert^2
    \ge
    \mu\bigl(f(x)-f^\star\bigr).
\]
В дальнейшем сильная выпуклость используется для нижней квадратичной
оценки и PL-неравенства, а гладкость — для верхней квадратичной
оценки.

\section{Ошибка по значению функции}

Положим
\[
    V_f(x)=f(x)-f^\star.
\]
При сильной выпуклости эта функция положительно определена
относительно \(x^\star\), а вдоль градиентного потока
\[
\begin{aligned}
    \dot V_f(x)
    &=
    \nabla f(x)^{\mathsf T}\dot x
    \\
    &=
    -\lVert\nabla f(x)\rVert^2.
\end{aligned}
\]
По PL-неравенству
\[
    \dot V_f\le-2\mu V_f,
\]
откуда
\[
    f(x(t))-f^\star
    \le
    e^{-2\mu(t-t_0)}
    \bigl(f(x(t_0))-f^\star\bigr).
\]
Значение целевой функции сходится к минимуму экспоненциально.

Чтобы получить оценку состояния по общему квадратичному критерию,
используем
\[
    \dot V_f
    \le
    -2\mu V_f
    \le
    -\mu^2\lVert x-x^\star\rVert^2.
\]
Можно выбрать
\[
    c_1=\frac\mu2,
    \qquad
    c_2=\frac L2,
    \qquad
    c_3=\mu^2.
\]
Следовательно,
\[
    \lVert x(t)-x^\star\rVert
    \le
    \sqrt{\frac L\mu}
    \exp\left(
        -\frac{\mu^2}{L}(t-t_0)
    \right)
    \lVert x(t_0)-x^\star\rVert.
\]
В терминах числа обусловленности
\[
    \kappa=\frac L\mu
\]
эта оценка принимает вид
\[
    \lVert x(t)-x^\star\rVert
    \le
    \sqrt\kappa
    \exp\left(
        -\frac\mu\kappa(t-t_0)
    \right)
    \lVert x(t_0)-x^\star\rVert.
\]
Чем больше \(\kappa\), тем слабее полученная граница. Это свойство
данного доказательства через \(V_f\), а не точная характеристика
самого градиентного потока.

\section{Расстояние и норма градиента}

Более прямой выбор функции Ляпунова равен
\[
    V_x(x)
    =
    \frac12\lVert x-x^\star\rVert^2.
\]
Тогда
\[
    \dot V_x
    =
    -\bigl\langle
        x-x^\star,
        \nabla f(x)
    \bigr\rangle.
\]
Сильная монотонность градиента даёт
\[
    \bigl\langle
        \nabla f(x)-\nabla f(x^\star),
        x-x^\star
    \bigr\rangle
    \ge
    \mu\lVert x-x^\star\rVert^2.
\]
Поскольку \(\nabla f(x^\star)=0\), получаем
\[
    \dot V_x
    \le
    -\mu\lVert x-x^\star\rVert^2
    =
    -2\mu V_x.
\]
Следовательно,
\[
    \lVert x(t)-x^\star\rVert
    \le
    e^{-\mu(t-t_0)}
    \lVert x(t_0)-x^\star\rVert.
\]
Эта оценка сильнее лекционного вывода через \(f-f^\star\) и не
требует \(L\)-гладкости для самой оценки, если траектория существует.
Гладкость нужна только в предыдущем доказательстве, где используется
верхняя квадратичная оценка на \(f-f^\star\).

Если \(f\in C^2\), можно также взять
\[
    V_{\nabla}(x)
    =
    \frac12\lVert\nabla f(x)\rVert^2.
\]
Тогда
\[
\begin{aligned}
    \dot V_{\nabla}
    &=
    \nabla f(x)^{\mathsf T}
    \nabla^2f(x)\dot x
    \\
    &=
    -\nabla f(x)^{\mathsf T}
    \nabla^2f(x)\nabla f(x)
    \\
    &\le
    -\mu\lVert\nabla f(x)\rVert^2
    =
    -2\mu V_{\nabla}.
\end{aligned}
\]
Поэтому
\[
    \lVert\nabla f(x(t))\rVert
    \le
    e^{-\mu(t-t_0)}
    \lVert\nabla f(x(t_0))\rVert.
\]
При \(L\)-гладкости и сильной выпуклости
\[
    \mu\lVert x-x^\star\rVert
    \le
    \lVert\nabla f(x)\rVert
    \le
    L\lVert x-x^\star\rVert,
\]
откуда следует ещё одна оценка состояния:
\[
    \lVert x(t)-x^\star\rVert
    \le
    \frac L\mu
    e^{-\mu(t-t_0)}
    \lVert x(t_0)-x^\star\rVert.
\]

Функции \(V_f\), \(V_x\) и \(V_{\nabla}\) контролируют разные
характеристики: ошибку по значению, расстояние до решения и норму
градиента. Выбор зависит от того, какую величину требуется оценить и
какие предположения доступны.

\section{PL-условие без сильной выпуклости}

Пусть функция удовлетворяет PL-неравенству
\[
    \frac12\lVert\nabla f(x)\rVert^2
    \ge
    \mu\bigl(f(x)-f^\star\bigr),
\]
но не обязательно является сильно выпуклой. Для
\[
    V_f(x)=f(x)-f^\star
\]
по-прежнему выполняется
\[
    \dot V_f
    =
    -\lVert\nabla f(x)\rVert^2
    \le
    -2\mu V_f,
\]
и потому
\[
    f(x(t))-f^\star
    \le
    e^{-2\mu(t-t_0)}
    \bigl(f(x(t_0))-f^\star\bigr).
\]

Однако PL-свойство само по себе не даёт двусторонней связи между
ошибкой по значению и расстоянием до единственной точки. Множество
глобальных минимумов может содержать несколько элементов, поэтому в
общем случае естественно исследовать устойчивость множества
\[
    \mathcal X^\star
    =
    \operatorname*{argmin}f.
\]
Для перехода от ошибки по функции к расстоянию
\[
    \operatorname{dist}(x,\mathcal X^\star)
\]
нужно дополнительное условие роста или регулярности.

Из PL-неравенства следует, что всякая стационарная точка является глобальным минимумом.

\section{Связь с дискретным методом}

Градиентный поток является непрерывным аналогом итерации
\[
    x^{k+1}
    =
    x^k-\eta\nabla f(x^k).
\]
Для \(\mu\)-сильно выпуклой \(L\)-гладкой функции и
\[
    \eta=\frac1L
\]
дискретный метод удовлетворяет оценке
\[
    f(x^k)-f^\star
    \le
    \left(1-\frac\mu L\right)^k
    \bigl(f(x^0)-f^\star\bigr).
\]
В непрерывном времени соответствующая функциональная оценка имеет
вид
\[
    f(x(t))-f^\star
    \le
    e^{-2\mu t}
    \bigl(f(x(0))-f^\star\bigr).
\]
Эти скорости нельзя сравнивать только по коэффициентам: один
дискретный шаг имеет вычислительную стоимость, а масштаб времени
непрерывной системы изменяется при умножении правой части на
положительную константу.

Практическое доказательство экспоненциальной устойчивости состоит из
четырёх действий: определить равновесие, выбрать функцию Ляпунова,
связать её с нужной мерой ошибки и оценить производную. Неравенство
\[
    \dot V\le-\gamma V
\]
непосредственно даёт экспоненциальное убывание выбранной величины.
Чтобы получить экспоненциальную устойчивость состояния, дополнительно
нужна количественная связь между \(V\) и расстоянием до равновесия.
